\chapter{Dispersion et séries regroupées en classes}
\section*{Exercice 19.1 -- Calculer une dispersion}
Pour la série $2,4,4,6$, calculer moyenne, variance et écart-type. Contrôler la variance de deux façons.

\section*{Exercice 19.2 -- Changer d'unité}
Une série de températures a une moyenne de $18\degre$C et un écart-type de $3\degre$C. La conversion est $F=1{,}8C+32$. Calculer les deux indicateurs en degrés Fahrenheit et justifier l'absence de $32$ dans le second.

\section*{Exercice 19.3 -- Classes et histogramme}
Les classes $[0;10[$, $[10;20[$, $[20;30]$ ont pour effectifs $4,8,8$. Tracer un histogramme d'amplitudes égales et le polygone des fréquences cumulées. Estimer moyenne et médiane sous l'hypothèse de répartition uniforme dans chaque classe.

\section*{Exercice 19.4 -- Ce que le regroupement perd}
Avec les mêmes classes et effectifs, les moyennes exactes de classe sont $4$, $14$ et $26$. Déterminer la moyenne exacte globale. Peut-on en déduire l'écart-type exact ? Construire une explication.

\section*{Exercice 19.5 -- Choisir les indicateurs}
Comparer les séries $A=(10,11,12,13,14)$ et $B=(10,11,12,13,100)$. Calculer médiane, quartiles, moyenne et étendue. Expliquer pourquoi un seul indicateur serait trompeur.

\section*{Exercice 19.6 -- Décomposer une variance (facultatif)}
Deux groupes de même effectif ont des moyennes $2$ et $8$ et des variances $1$ et $4$. Calculer la moyenne et la variance de la réunion. Pourquoi la moyenne des variances ne convient-elle pas ?
